\scanpage{26}
% archive sheet number: 25
% L1--3
Les fibres de ${\unsure{\mathscr P|\mathcal P}}^0_{X/S}$ sont connexes, et propres si celles de $X$ sont \uline{simples}. Soit \deleted{{\unsure{schéma}}} $X-S$ l’{\unsure{ens.|ouv.}} des pts de $X$ qui sont {\unsure{non}} simples sur $Y$. On va définir un morphisme
% L4
\[
\boxed{X-S\longrightarrow\mathscr P^1_{X/S}}.
\]
% L5--8
où l’exposant $1$ indique qu’on est dans la composante d’augmentation $1$. {\unsure{Je reviens|Il revient}} {\unsure{à en donner}} \uline{fonctoriellement} \inserted{pour le changement de base} une application
\deleted{{\unsure{application}}}
\[
{\unsure{\Gamma|\mathscr T}}(X-S/Y)\longrightarrow\mathscr P^1(X/S)
\]
% L9--13
or on note qu’une section de $X-S$ sur $Y$ \uline{définit un diviseur sur} $X-S$ {\unsure{(transverse à $X/Y$)}} d’où un fibré inversible, dont on constate que le \deleted{{\unsure{diviseur}}} \inserted{degré} est $+1$ [car on peut supposer ${\unsure{Y|S}}=\operatorname{Spec}k$, {\unsure{le degré clair}}].
% margin note written slanted, rising toward the text
\margin{{\unsure{vérifier|référence}}??}

% L14--23
II) Partons d’un $X$ \uline{propre} et \uline{simple} sur ${\unsure{Y|S}}$, dont les fibres sont de dim $1$, \deleted{{\unsure{Soit}} $X_i$} {\unsure{(et abs! connexes)}} (donc abs! irréductibles). Soit $S$ un sous-schéma \inserted{fermé} de \deleted{$X$ plat sur $Y$ et} ne contenant aucune fibre de $X/{\unsure{Y|S}}$, donc \uline{fini} sur ${\unsure{Y|S}}$, et plat sur $Y$. Soit ${\unsure{\Xi|\mathscr I|\mathcal I}}$ le faisceau d’idéaux définissant $Y$. Alors on prouve facilement que $\Xi$ est \uline{localement libre} [c’est purement local : un faisceau coh. sur $X$ plat rel. à $Y$ et qui {\unsure{est inversible|est un module}} dans fibres \inserted{{\unsure{est}}} est plat sur $X$.---]
% margin note written slanted, read after rotating clockwise
\margin{{\unsure{fibres abs! connexes est une condition}} {\unsure{peut-être}} \uline{{\unsure{vide}}}.}
